\documentclass{article}
\usepackage{geometry}
\geometry{margin=1.5cm, vmargin={0pt,1cm}}
\setlength{\topmargin}{-1cm}
\setlength{\headheight}{12.5pt}
\setlength{\paperheight}{29.7cm}
\setlength{\textheight}{25.3cm}

% useful packages.
% \usepackage[UTF8]{ctex}
\usepackage{fancyhdr}
\usepackage{layout}
\usepackage{tikz}
\usetikzlibrary{arrows.meta}

% import common macros
% % % % % % % % % % % % % % % % % % % % % % % % % % % % % % % %
% Common Macros for LaTeX
% % % % % % % % % % % % % % % % % % % % % % % % % % % % % % % %

% Useful packages
\usepackage{amsfonts}
\usepackage{amsmath}
\usepackage{amssymb}
\usepackage{amsthm}
\usepackage{enumerate}
\usepackage{graphicx}
\usepackage{multicol}
\usepackage[ruled,linesnumbered]{algorithm2e}

% Math operators and common symbols
\newcommand{\dif}{\mathrm{d}}
\newcommand{\innerProd}[1]{\left\langle #1 \right\rangle}
\newcommand{\difFrac}[2]{\frac{\dif #1}{\dif #2}}
\newcommand{\pdfFrac}[2]{\frac{\partial #1}{\partial #2}}
\newcommand{\T}{\mathrm{T}}
\newcommand{\norm}[1]{\left\| #1 \right\|}
\newcommand{\abs}[1]{\left| #1 \right|}

% Vector and Matrix shortcuts
\newcommand{\vecTwo}[2]{
  \begin{bmatrix}
    #1 \\
    #2
\end{bmatrix}}
\newcommand{\vecThree}[3]{
  \begin{bmatrix}
    #1 \\
    #2 \\
    #3
\end{bmatrix}}
\newcommand{\vecTwoH}[2]{
  \begin{bmatrix}
    #1 & #2
  \end{bmatrix}
}
\newcommand{\vecThreeH}[3]{
  \begin{bmatrix}
    #1 & #2 & #3
  \end{bmatrix}
}

% Common fields
\newcommand{\R}{\mathbb{R}}
\newcommand{\C}{\mathbb{C}}
\newcommand{\Z}{\mathbb{Z}}
\newcommand{\N}{\mathbb{N}}

% Solutions and proofs
\newenvironment{solution}{
\begin{proof}[解]}{
\end{proof}}

% Utility to get environment variables (requires lualatex)
\newcommand{\getEnv}[2]{\directlua{tex.print(os.getenv("#1") or "#2")}}


% Name and uID
\newcommand{\name}{\getEnv{NAME}{NAME}}
\newcommand{\uid}{\getEnv{UID}{UID}}
\newcommand{\course}{Complex Analysis}
\newcommand{\hwNum}{8}

\begin{document}

\pagestyle{fancy}
\fancyhead{}
\lhead{\zhstr{\name} (\uid)}
\chead{\course\ \#\hwNum}
\rhead{\today}

\section{Exercise 4.7}

Denote $S^1 = \{e^{it} : t \in [0, 2\pi]\}$ and consider a closed curve
$\Gamma$ as the image of $S^1$ under the function
\[
  \gamma(t) = z + r(t)\,e^{i\theta(t)},
\]
where $\theta$ is the angle of $\gamma(t)$ in the polar coordinate system
whose pole is at $z$.  Show that
\[
  w(\Gamma, z) = \frac{1}{2\pi}\bigl(\theta(2\pi) - \theta(0)\bigr).
\]
According to this result, what is the meaning of the winding number
$w(\Gamma, z)$?

\begin{solution}
  Since $\gamma(t) - z = r(t)\,e^{i\theta(t)}$, we compute
  \[
    \frac{d\gamma}{\gamma - z}
    = \frac{d\bigl(r\,e^{i\theta}\bigr)}{r\,e^{i\theta}}
    = \frac{r' + ir\theta'}{r}\,dt
    = \frac{r'}{r}\,dt + i\,\theta'\,dt.
  \]
  By Definition~4.2,
  \[
    w(\Gamma, z)
    = \frac{1}{2\pi i}\int_0^{2\pi}
      \!\left(\frac{r'}{r} + i\theta'\right)dt
    = \frac{1}{2\pi i}\Bigl[\ln r(t)\Bigr]_0^{2\pi}
      + \frac{1}{2\pi}\Bigl[\theta(t)\Bigr]_0^{2\pi}.
  \]
  Since $\Gamma$ is closed, $r(2\pi) = r(0)$, so
  $\bigl[\ln r(t)\bigr]_0^{2\pi} = 0$.
  Therefore
  \[
    w(\Gamma, z) = \frac{\theta(2\pi) - \theta(0)}{2\pi}.
  \]

  This means the winding number $w(\Gamma, z)$ counts the number of
  complete revolutions that $\gamma(t)$ makes around $z$ as $t$ traverses
  $[0, 2\pi]$.\qedhere
\end{solution}

\section{Exercise 4.11}

\textbf{Lemma 4.10.}
If a particle crosses an oriented closed curve a finite number of times
as it moves continuously from $x$ to infinity, the winding number
$w_\gamma(x)$ is equal to the number of negative crossings minus the
number of positive crossings.

Prove Lemma 4.10 by induction.

\begin{solution}
  We proceed by induction on the number $n$ of crossings.

  \textbf{Base case ($n = 0$).}\;
  The particle moves from $x$ to infinity without crossing $\Gamma$,
  so $x$ and infinity lie in the same connected component of
  $\C \setminus \Gamma$.  The unbounded component contains infinity,
  so $x$ is in the unbounded component.  By Lemma~4.4,
  $w(\Gamma, x) = 0 = 0 - 0$.\quad$\checkmark$

  \textbf{Inductive step.}\;
  Assume the result holds for $n$ crossings.
  Consider a path from $x$ to infinity with $n + 1$ crossings,
  and let $t_c$ be the time of the first crossing.

  Before $t_c$, the particle has not crossed $\Gamma$, so $x$ and
  $p(t_c^-)$ lie in the same component of $\C \setminus \Gamma$.
  By Lemma~4.3, $w(\Gamma, x) = w(\Gamma, p(t_c^-))$.

  After $t_c$, the particle makes $n$ remaining crossings.
  By the inductive hypothesis,
  $w(\Gamma, p(t_c^+)) = N_r^- - N_r^+$,
  where $N_r^-$ and $N_r^+$ are the remaining negative and positive
  crossings.

  By Lemma~4.9:
  \begin{itemize}
    \item \textit{Positive crossing} ($w$ increases by~$1$):
      $w(\Gamma, p(t_c^+)) = w(\Gamma, p(t_c^-)) + 1$, so
      \[
        w(\Gamma, x) = w(\Gamma, p(t_c^+)) - 1
        = (N_r^- - N_r^+) - 1
        = N_r^- - (N_r^+ + 1)
        = N^- - N^+.
      \]
    \item \textit{Negative crossing} ($w$ decreases by~$1$):
      $w(\Gamma, p(t_c^+)) = w(\Gamma, p(t_c^-)) - 1$, so
      \[
        w(\Gamma, x) = w(\Gamma, p(t_c^+)) + 1
        = (N_r^- - N_r^+) + 1
        = (N_r^- + 1) - N_r^+
        = N^- - N^+.
      \]
  \end{itemize}
  In both cases $w(\Gamma, x) = N^- - N^+$.\qedhere
\end{solution}

\section{Exercise 4.13}

\textbf{Theorem 4.12.}
Suppose the cycle decomposition of an oriented closed curve $\Gamma$
as in Lemma 3.68 consists of a finite number of Jordan curves.
Then for each $x \notin \Gamma$, its winding number $w(\Gamma, x)$ is
the number of counterclockwise Jordan curves containing $x$ minus that
of clockwise Jordan curves containing $x$.

Prove Theorem 4.12.

\begin{solution}
  By the Seifert decomposition (Lemma~3.68), $\Gamma$ decomposes into
  a finite number of Jordan curves $\Gamma_1, \ldots, \Gamma_k$
  with induced orientations.
  By additivity of the line integral,
  \[
    w(\Gamma, x)
    = \frac{1}{2\pi i}\int_\Gamma \frac{d\gamma}{\gamma - x}
    = \sum_{j=1}^{k} w(\Gamma_j, x).
  \]

  For each $\Gamma_j$, since $x \notin \Gamma \supseteq \Gamma_j$,
  exactly one of the following holds by Corollaries~4.32 and~4.33:
  \begin{itemize}
    \item $\Gamma_j$ positively oriented, $x \in \mathrm{BCJ}(\Gamma_j)$:
      $w(\Gamma_j, x) = 1$.
    \item $\Gamma_j$ negatively oriented, $x \in \mathrm{BCJ}(\Gamma_j)$:
      $w(\Gamma_j, x) = -1$.
    \item $x \notin \mathrm{BCJ}(\Gamma_j)$:
      $w(\Gamma_j, x) = 0$.
  \end{itemize}
  Therefore
  \[
    w(\Gamma, x)
    = \#\bigl\{j : \text{$\Gamma_j$ counterclockwise, }
      x \in \mathrm{BCJ}(\Gamma_j)\bigr\}
    - \#\bigl\{j : \text{$\Gamma_j$ clockwise, }
      x \in \mathrm{BCJ}(\Gamma_j)\bigr\}.\qedhere
  \]
\end{solution}

\section{Exercise 4.15}

\textbf{Lemma 4.14.}
Any closed curve $\Gamma$ that satisfies $0 \notin \Gamma$ and
$w(\Gamma, 0) = \nu \ne 0$ can be continuously deformed into the
archetypal loop $\hat{\gamma}_\nu = e^{i\nu\phi}$ with
$\phi \in [0, 2\pi]$, and vice versa.

For the proof of Lemma 4.14, give another free homotopy in which the
modulus and the argument change simultaneously from $0$ to $1$.

\begin{solution}
  Parametrize $\Gamma$ as $\gamma(\phi) = r(\phi)\,e^{i\Phi(\phi)}$
  with $\phi \in [0, 2\pi]$, where
  $\Phi(2\pi) - \Phi(0) = 2\pi\nu$.
  The archetypal loop is
  $\hat{\gamma}_\nu(\phi) = e^{i\nu\phi}$.

  Define the free homotopy
  \[
    H(s, t)
    := \bigl[(1-t)\,r(2\pi s) + t\bigr]
    \cdot \exp\bigl\{i\bigl[(1-t)\,\Phi(2\pi s)
    + t \cdot 2\pi\nu s\bigr]\bigr\}
  \]
  for $s \in [0, 1]$ and $t \in [0, 1]$.

  \textbf{Endpoints.}\;
  $H(s, 0) = r(2\pi s)\,e^{i\Phi(2\pi s)} = \gamma(2\pi s)$
  and
  $H(s, 1) = 1 \cdot e^{i\,2\pi\nu s}
  = \hat{\gamma}_\nu(2\pi s)$.

  \textbf{Closedness.}\;
  Since $\gamma$ is closed, $r(0) = r(2\pi)$ and
  $\Phi(2\pi) = \Phi(0) + 2\pi\nu$.
  Then
  \begin{align*}
    H(1, t)
    &= \bigl[(1-t)\,r(0) + t\bigr]\,
       e^{i\bigl[(1-t)(\Phi(0)+2\pi\nu)+2\pi\nu t\bigr]} \\
    &= \bigl[(1-t)\,r(0) + t\bigr]\,
       e^{i(1-t)\Phi(0)}\,e^{i\,2\pi\nu} \\
    &= H(0, t) \cdot e^{i\,2\pi\nu}
    = H(0, t),
  \end{align*}
  where the last step uses $\nu \in \Z$ (Lemma~4.3).

  \textbf{Non-vanishing.}\;
  Since $0 \notin \Gamma$, $r(\phi) > 0$ for all $\phi$,
  so $(1-t)\,r(2\pi s) + t \ge \min(r, 1) > 0$.

  Both the modulus and the argument vary simultaneously
  from their values on $\Gamma$ to those on $\hat{\gamma}_\nu$
  as $t$ goes from $0$ to $1$.\qedhere
\end{solution}

\section{Exercise 4.36 (Alternative proof of the fundamental theorem of algebra)}

Prove Theorem 3.47 (the fundamental theorem of algebra) via the argument
principle in Lemma 4.34.

\textbf{Lemma 4.34.}
Suppose a polynomial $f$ has no roots on a positively oriented Jordan
curve $\Gamma$.  Then $w(f(\Gamma), 0)$ equals the number of roots of
$f$ in $\mathrm{BCJ}(\Gamma)$, counted with their algebraic multiplicity.

\begin{solution}
  \textbf{Theorem 3.47 (Fundamental theorem of algebra).}
  Every polynomial of degree $n \ge 1$ with coefficients in $\C$
  has at least one root in $\C$.

  \medskip

  Let $f(z) = a_n z^n + a_{n-1}z^{n-1} + \cdots + a_0$
  with $n \ge 1$ and $a_n \ne 0$.
  Write $f(z) = a_n z^n\,g(z)$ where
  \[
    g(z) = 1 + \frac{a_{n-1}}{a_n z} + \cdots + \frac{a_0}{a_n z^n}.
  \]
  Since $g(z) \to 1$ as $|z| \to \infty$, there exists $R > 0$
  such that $|g(z) - 1| < 1$ for $|z| \ge R$.
  In particular, $g(z) \ne 0$ for $|z| \ge R$, and hence $f$
  has no roots on $|z| = R$.

  Let $\Gamma_R$ be the positively oriented circle $|z| = R$.
  By Lemma~4.34,
  \[
    w(f(\Gamma_R),\, 0)
    = \text{number of roots of } f \text{ in } |z| < R
      \text{ (with multiplicity)}.
  \]
  Since $\frac{f'}{f} = \frac{n}{z} + \frac{g'}{g}$ and
  $g(\Gamma_R)$ lies in the disk $|w - 1| < 1$ (which does not
  contain~$0$), we have $w(g(\Gamma_R), 0) = 0$.
  Therefore
  \[
    w(f(\Gamma_R),\, 0)
    = \frac{1}{2\pi i}\oint_{\Gamma_R}\!\frac{n}{z}\,dz
    + \frac{1}{2\pi i}\oint_{\Gamma_R}\!\frac{g'(z)}{g(z)}\,dz
    = n + 0 = n \ge 1.
  \]
  Hence $f$ has at least one root in $|z| < R \subset \C$.\qedhere
\end{solution}

\section{Exercise 4.38}

\textbf{Lemma 4.37.}
Let $f, g : \Omega \to \C$ be analytic functions where
$\Omega \subset \C$ is a simply connected region.  Suppose
$\Gamma \subset \Omega$ is an oriented closed curve such that
$f(z)\,g(z) \ne 0$ for all $z \in \Gamma$.  Then we have
\[
  w\bigl((fg)(\Gamma),\, 0\bigr)
  = w\bigl(f(\Gamma),\, 0\bigr)
  + w\bigl(g(\Gamma),\, 0\bigr).
\]

Prove Lemma 4.37.

\begin{solution}
  By Definition~4.2,
  \[
    w\bigl((fg)(\Gamma),\, 0\bigr)
    = \frac{1}{2\pi i}\oint_\Gamma
      \frac{(fg)'(\gamma)}{(fg)(\gamma)}\,d\gamma.
  \]
  Since $(fg)' = f'g + fg'$ and
  $f(\gamma)\,g(\gamma) \ne 0$ for all $\gamma \in \Gamma$,
  \[
    \frac{(fg)'}{fg} = \frac{f'}{f} + \frac{g'}{g}.
  \]
  Therefore
  \[
    w\bigl((fg)(\Gamma),\, 0\bigr)
    = \frac{1}{2\pi i}\oint_\Gamma \frac{f'(\gamma)}{f(\gamma)}\,d\gamma
    + \frac{1}{2\pi i}\oint_\Gamma \frac{g'(\gamma)}{g(\gamma)}\,d\gamma
    = w\bigl(f(\Gamma),\, 0\bigr) + w\bigl(g(\Gamma),\, 0\bigr).\qedhere
  \]
\end{solution}

\section{Exercise 4.41}

\textbf{Corollary 4.40.}
Let $\Omega$ be a domain and $\Gamma : [0, 1] \to \Omega$ be a
positively oriented Jordan curve inside $\Omega$.  For analytic
functions $f, h : \Omega \to \C$ and a point $p \notin f(\Gamma)$, we
have
\[
  \frac{1}{2\pi i} \oint_\Gamma h(z)\,\frac{f'(z)}{f(z) - p}\,dz
  = \sum_{z_j \in f^{-1}(p)} h(z_j)\,m_j,
\]
where $m_j$ is the multiplicity of the root $z_j$ of $f(x) = p$.

Prove Corollary 4.40.

\begin{solution}
  Let $z_1, \ldots, z_k$ be the distinct roots of $f(z) = p$
  inside $\Gamma$ with multiplicities $m_1, \ldots, m_k$.
  By Lemma~4.22, near each $z_j$,
  \[
    f(z) - p = (z - z_j)^{m_j}\,\phi_j(z),
  \]
  where $\phi_j$ is analytic and $\phi_j(z_j) \ne 0$.
  Hence
  \[
    \frac{f'(z)}{f(z) - p}
    = \frac{m_j}{z - z_j} + \frac{\phi_j'(z)}{\phi_j(z)},
  \]
  and the second term is analytic near $z_j$.

  Deform $\Gamma$ into small positively oriented circles
  $\Gamma_j$ around each $z_j$, each lying in a neighborhood
  where the above factorization holds and containing no other
  roots.  By Cauchy's theorem,
  \[
    \frac{1}{2\pi i}\oint_\Gamma h(z)\,\frac{f'(z)}{f(z)-p}\,dz
    = \sum_{j=1}^{k}
      \frac{1}{2\pi i}\oint_{\Gamma_j}
      h(z)\,\frac{f'(z)}{f(z)-p}\,dz.
  \]
  On each $\Gamma_j$,
  \[
    h(z)\,\frac{f'(z)}{f(z) - p}
    = \frac{m_j\,h(z)}{z - z_j}
    + h(z)\,\frac{\phi_j'(z)}{\phi_j(z)}.
  \]
  Since $h\,\phi_j'/\phi_j$ is analytic near $z_j$,
  its contour integral vanishes.
  By the Cauchy integral formula,
  $\frac{1}{2\pi i}\oint_{\Gamma_j} \frac{m_j\,h(z)}{z - z_j}\,dz
  = m_j\,h(z_j)$.
  Therefore
  \[
    \frac{1}{2\pi i}\oint_\Gamma
    h(z)\,\frac{f'(z)}{f(z) - p}\,dz
    = \sum_{j=1}^{k} m_j\,h(z_j).\qedhere
  \]
\end{solution}

\section{Exercise 4.47}

\textbf{Theorem 4.46 (Darboux--Picard).}
Suppose $f : \Omega \to \C$ is an analytic function where $\Omega$ is
a domain that contains a Jordan curve $\Gamma : [0, 1] \to \Omega$.
If the restriction $f|_\Gamma$ is injective, then so is
$f|_{\mathrm{BCJ}(\Gamma)}$.

Give an alternative proof of Theorem 4.46 by considering points in
$\mathrm{BCJ}(\Gamma)$.

\begin{solution}
  Since $f|_\Gamma$ is injective and $\Gamma$ is a Jordan curve,
  $f(\Gamma)$ is a Jordan curve.

  Suppose for contradiction that
  $f(z_1) = f(z_2) = p$ for distinct $z_1, z_2 \in \mathrm{BCJ}(\Gamma)$.

  \textbf{Claim: $p \notin f(\Gamma)$.}\;
  If $p \in f(\Gamma)$, then by the open mapping theorem
  (Theorem~3.97), $f(\mathrm{BCJ}(\Gamma))$ is an open set
  containing $p$.
  Since $f(\Gamma)$ is a Jordan curve passing through $p$,
  this open set must contain points on both sides of $f(\Gamma)$
  near $p$.
  In particular, there exists
  $p' \in f(\mathrm{BCJ}(\Gamma)) \cap \mathrm{UCJ}(f(\Gamma))$,
  say $p' = f(z')$ with $z' \in \mathrm{BCJ}(\Gamma)$.
  By (4.10),
  \[
    w(f(\Gamma), p')
    = \sum_{z_j \in f^{-1}(p') \cap \mathrm{BCJ}(\Gamma)} m_j
    \ge m_{z'} \ge 1.
  \]
  But $p' \in \mathrm{UCJ}(f(\Gamma))$ implies
  $w(f(\Gamma), p') = 0$ by Corollary~4.32, a contradiction.

  \textbf{Deriving the contradiction.}\;
  Since $p \notin f(\Gamma)$, by (4.11),
  \[
    w(f(\Gamma), p)
    = \sum_{z_j \in f^{-1}(p) \cap \mathrm{BCJ}(\Gamma)} m_j
    \ge m_{z_1} + m_{z_2} \ge 2.
  \]
  But $f(\Gamma)$ is a Jordan curve and $p \notin f(\Gamma)$,
  so $w(f(\Gamma), p) \in \{-1, 0, 1\}$.
  This contradicts $w(f(\Gamma), p) \ge 2$.

  Hence $f$ is injective on $\mathrm{BCJ}(\Gamma)$.\qedhere
\end{solution}

\end{document}
